% ============ 第2章 §2.17 Pick插值问题(原书 pp.173-179) ============
\section{Pick 插值问题}\label{sec:2.17}
\index{插值问题!最小范数插值问题}


对 D 中相异的有限个点 $\{z_{j}\}_{j=1}^{n}$ 与 $\mathbb{C}$ 中有限个数据 $\{w_{j}\}_{j=1}^{n}$,求属于 $H^{\infty}$ 且满足
\begin{equation*}
f(z_{j}) = w_{j} \quad (1 \leq j \leq n)
\end{equation*}
的函数 f.事实上,对 $1 \leq j \leq n$,令
\begin{equation*}
B_{n}(z) = \prod_{j=1}^{n} \frac{z - z_{j}}{1 - \bar{z}_{j}z} \quad \text{且} \quad B_{nj}(z) = B_{n}(z)\frac{1 - \bar{z}_{j}z}{z - z_{j}}.
\end{equation*}
令
\begin{equation*}
f(z) = \sum_{j=1}^{n} w_{j}\frac{B_{nj}(z)}{B_{nj}(z_{j})},
\end{equation*}
则 $f(z_{j}) = w_{j}$($1 \leq j \leq n$).于是当对 f 没有进一步条件时,有限个点的情形插值问题很容易解决.应用上重要的(有限)插值问题是:何时能选取 $\|f\|_{\infty} \leq 1$.

定理~\ref{thm:2.17.1} 称为最小范数插值问题,将在 2.18 节中使用.定理~\ref{thm:2.17.7} 称为 Pick 定理 \cite{ch2:14}.我们用 Sarason 的新方法,即利用定理~\ref{thm:2.17.5} 证明 Pick 定理.该方法不用复函数论,而用算子论.

\begin{theorem}\label{thm:2.17.1}
设 $1 \leq p \leq \infty$,$\{z_{j}\}_{j=1}^{n} \subset D$ 且 $\{w_{j}\}_{j=1}^{n} \subset \mathbb{C}$.此时存在满足 $f(z_{j}) = w_{j}$($1 \leq j \leq n$)的最小范数的 $f \in H^{p}$.且
\begin{align*}
\min \{\|f\|_{p}; f \in H^{p}, f(z_{j}) = w_{j} (1 \leq j \leq n)\}&= \max \left\{\left| \sum_{j=1}^{n} \frac{w_{j}(1 - |z_{j}|^{2})}{B_{nj}(z_{j})} g(z_{j}) \right|; g \in H^{q}, \|g\|_{q} \leq 1 \right\}.
 \end{align*}
\end{theorem}

\begin{proof}
首先,
\begin{equation*}
F(z) = \sum_{j=1}^{n} w_{j}\frac{B_{nj}(z)}{B_{nj}(z_{j})},
\end{equation*}
则 $F \in H^{p}$.对 $g \in H^{p}$,令
\begin{equation*}
f(z) = F(z) - B_{n}(z)g(z),
\end{equation*}
则 $f(z_{j}) = w_{j}$($1 \leq j \leq n$).对 $\{w_{j}\}_{j=1}^{n}$ 的一般插值函数具有这种形式.令 k = $F/B_{n}$,则 $k \in L^{p}$,且
\begin{equation*}
\inf_{g \in H^{p}} \|F - B_{n}g\|_{p} = \inf_{g \in H^{p}} \|k - g\|_{p}.
\end{equation*}
由定理~\ref{thm:2.14.2},inf 可改为 min,故存在满足 $f(z_{j}) = w_{j}$($1 \leq j \leq n$)的最小范数的 $f \in H^{p}$.

因
\begin{equation*}
k(z) = \sum_{j=1}^{n} \frac{w_{j}}{B_{nj}(z_{j})}\frac{1 - \bar{z}_{j}z}{z - z_{j}},
\end{equation*}
对 $h \in H^{q}$($1/p + 1/q = 1$),令
\begin{equation*}
\phi(f) = \frac{1}{2\pi i}\int_{-\pi}^{\pi} k(z) h(z)\dd z = \sum_{j=1}^{n} \frac{w_{j}(1 - |z_{j}|^{2})}{B_{nj}(z_{j})} h(z_{j}),
\end{equation*}
则 $\phi \in (H^{q})^{*}$.由定理~\ref{thm:2.14.2},当 $p \neq 1$ 时
\begin{align*}
\min \{\|f\|_{p}; f \in H^{p}, f(z_{j}) = w_{j} (1 \leq j \leq n)\}&= \min \{\|k - g\|_{p}; g \in H^{p}\}  \\
\quad &= \max \left\{\left| \sum_{j=1}^{n} \frac{w_{j}(1 - |z_{j}|^{2})}{B_{nj}(z_{j})} h(z_{j}) \right|; h \in H^{q}, \|h\|_{q} \leq 1 \right\}.
\end{align*}
p = 1 时,因 $k \in C$,由定理~\ref{thm:2.14.2} 与定理~\ref{thm:2.14.4} 的 (3),上述等式成立.
\end{proof}

对 $f \in H^{\infty}$,设 $M_{f}$ 为 $H^{2}$ 上的乘法算子,则它是 $H^{2}$ 上的有界线性算子.若在 $H^{2}$ 上定义内积 $\langle\ ,\ \rangle$ 为
\begin{equation*}
\langle f, g \rangle = \frac{1}{2\pi}\int_{-\pi}^{\pi} f(e^{i\theta})\overline{g(e^{i\theta})}\dd \theta \quad (f, g \in H^{2}),
\end{equation*}
则 $H^{2}$ 成为 Hilbert 空间.对内函数 q,设 $N = H^{2} \ominus qH^{2}$,即 $qH^{2}$ 在 $H^{2}$ 中的正交补空间.设 P 为从 $H^{2}$ 到 N 的正交射影.定义 N 上的有界线性算子 $T_{f}$ 为 $T_{f} = PM_{f}|N$.$T_{f}^{*}$ 表示 $T_{f}$ 的共轭算子.定理~\ref{thm:2.17.5} 在 q 为一般内函数时由 Sarason \cite{ch2:15} 证明.

\begin{lemma}\label{lem:2.17.2}
当 $q = B_{n}$ 时,N 由 $\{(1 - \bar{z}_{j}z)^{-1}\}_{j=1}^{n}$ 生成.
\end{lemma}

\begin{proof}
当 $|a| < 1$ 时,由命题 2.9.4,若 $g \in H^{2}$ 则
\begin{equation*}
g(a) = \frac{1}{2\pi}\int_{-\pi}^{\pi} \frac{e^{it}}{e^{it} - a} g(e^{it})\dd t.
\end{equation*}
因
\begin{equation*}
\frac{e^{it}}{e^{it} - a} = \frac{1}{1 - ae^{-it}},
\end{equation*}
$\{(1 - \bar{z}_{j}z)^{-1}\}_{j=1}^{n}$ 与 $B_{n}H^{2}$ 正交.因 $B_{n}H^{2} = \{g \in H^{2}; g(z_{j}) = 0 \ (1 \leq j \leq n)\}$,N 是 n 维空间,故 N 由 $\{(1 - \bar{z}_{j}z)^{-1}\}_{j=1}^{n}$ 生成.
\end{proof}

\begin{lemma}\label{lem:2.17.3}
当 $q = B_{n}$ 时,若 $f \in H^{\infty}$,则
\begin{equation*}
T_{f}^{*}\left(\frac{1}{1 - \bar{z}_{j}z}\right) = \frac{\overline{f(z_{j})}}{1 - \bar{z}_{j}z} \quad (1 \leq j \leq n).
\end{equation*}
\end{lemma}

\begin{proof}
令 $k_{j} = (1 - \bar{z}_{j}z)^{-1}$($1 \leq j \leq n$),则对任意 $g \in N$,
\begin{align*}
\langle T_{f}^{*}k_{j}, g \rangle & = \langle k_{j}, fg \rangle = \overline{\langle fg, k_{j} \rangle} = \overline{f(z_{j})g(z_{j})} \\
& = \langle \overline{f(z_{j})}k_{j}, g \rangle,
\end{align*}
故 $T_{f}^{*}k_{j} = \overline{f(z_{j})}k_{j}$.
\end{proof}

\begin{lemma}\label{lem:2.17.4}
设 $q = B_{n}$.此时对任意 $f \in H^{\infty}$,$T_{z}T_{f} = T_{f}T_{z}$.而且若 A 是 N 上的有界线性算子且 $T_{z}A = AT_{z}$,则对任意 $h \in N$ 有 $T_{h}A = AT_{h}$.
\end{lemma}

\begin{proof}
对任意 $f, g \in H^{\infty}$,
\begin{align*}
T_{f}T_{g} & = PM_{f}PM_{g}|N = PM_{f}(P + I - P)M_{g}|N \\
& = PM_{fg}|N = T_{fg},
\end{align*}
这里用到 $(I-P)H^{2} = B_{n}H^{2}$ 故 $M_{f}(I-P)M_{g}N \subset B_{n}H^{2}$.这证明了前半部分.又当 $T_{z}A = AT_{z}$ 时,若 f 是 z 的多项式则 $T_{f}A = AT_{f}$.若 $h \in N$,由引理~\ref{lem:2.17.2},h 可由 z 的多项式一致逼近.利用这些事实与 $\|T_{f}\| \leq \|f\|_{\infty}$,可以证明 $T_{h}A = AT_{h}$.
\end{proof}

\begin{theorem}[Sarason]\label{thm:2.17.5}
设 $q = B_{n}$.若 A 是 N 上的有界线性算子且 $T_{z}A = AT_{z}$,则存在 $f \in H^{\infty}$,使得 $A = T_{f}$ 且 $\|A\| = \|f\|_{\infty}$.
\end{theorem}

\begin{proof}
由引理~\ref{lem:2.17.2},$N \subset H^{\infty}$,令 F = AP1,则 $F \in H^{\infty}$.由引理~\ref{lem:2.17.4},设 h, g $\in N$,则 $T_{h}T_{F} = T_{F}T_{h}$,故
\begin{align*}
\langle Ah, g \rangle & = \langle AT_{h}P1, g \rangle = \langle T_{h}AP1, g \rangle \\
& = \langle T_{h}T_{F}P1, g \rangle = \langle T_{F}T_{h}P1, g \rangle \\
& = \langle T_{F}h, g \rangle .
\end{align*}
因 $h, g \in N$ 任意,$A = T_{F}$,且 $\|A\| \leq \|F\|_{\infty}$.

若 $F - G \in B_{n}H^{\infty}$ 则 $T_{F} = T_{G}$,故 $\|A\| \leq \|F + B_{n}H^{\infty}\|$.若能证明 $\|A\| = \|F + B_{n}H^{\infty}\|$,则由定理~\ref{thm:2.14.2} 可证明存在 f 使得 $A = T_{f}$ 且 $\|A\| = \|f\|_{\infty}$.设 $h \in H^{2}$,$k \in B_{n}\overline{H}_{0}^{2}$,则
\begin{align*}
\langle Fh, k \rangle & = \langle FPh, k \rangle + \langle F(I - P)h, k \rangle \\
& = \langle FPh, k \rangle = \langle PFPh, k \rangle .
\end{align*}
于是令 $k = B_{n}\bar{g}$,则
\begin{equation*}
\left| \int_{-\pi}^{\pi} \bar{B}_{n}Fhg \dd \theta / 2\pi \right| \leq \|T_{F}\|\|h\|_{2}\|g\|_{2}.
\end{equation*}
任意 $\ell \in e^{i\theta}H^{1}$ 由定理~\ref{thm:2.8.3} 可因式分解为 $\ell = QL$,令 $h = L^{1/2}$ 且 $g = QL^{1/2}$,则 $\|\ell\|_{1} = \|h\|_{2}\|g\|_{2}$.于是
\begin{equation*}
\left| \int_{-\pi}^{\pi} \bar{B}_{n}F\ell \dd \theta / 2\pi \right| \leq \|T_{F}\|\|\ell\|_{1},
\end{equation*}
由定理~\ref{thm:2.14.2},
\begin{equation*}
\|\bar{B}_{n}F + H^{\infty}\| \leq \|T_{F}\|,
\end{equation*}
故可证 $\|T_{F}\| = \|F + B_{n}H^{\infty}\|$.
\end{proof}

\begin{lemma}\label{lem:2.17.6}
设 $q = B_{n}$.若 $B(1 - \bar{z}_{j}z)^{-1} = \bar{w}_{j}(1 - \bar{z}_{j}z)^{-1}$($1 \leq j \leq n$),则 B 可扩张为 N 上的有界线性算子 $\tilde{B}$,且 $\|\tilde{B}\| \leq 1$ 的充要条件是
\begin{equation*}
\sum_{j,k=1}^{n} \frac{1 - w_{j}\bar{w}_{k}}{1 - z_{j}\bar{z}_{k}} \alpha_{j}\bar{\alpha}_{k} \geq 0 \quad (\{\alpha_{j}\}_{j=1}^{n} \subset \mathbb{C}).
\end{equation*}
\end{lemma}

\begin{proof}
由引理~\ref{lem:2.17.2},若 $h \in N$,则 $h = \sum_{j=1}^{n} \alpha_{j}(1 - \bar{z}_{j}z)^{-1}$,令
\begin{equation*}
\tilde{B}h = \sum_{j=1}^{n} \alpha_{j}\bar{w}_{j}(1 - \bar{z}_{j}z)^{-1},
\end{equation*}
则 $\tilde{B}$ 是 N 上的有界线性算子.$\|h\|_{2}^{2} = \sum_{j,k}\alpha_{j}\bar{\alpha}_{k}(1 - \bar{z}_{j}z_{k})^{-1}$,且
\begin{equation*}
\|\tilde{B}h\|_{2}^{2} = \sum_{j,k}\alpha_{j}\bar{\alpha}_{k}\bar{w}_{j}w_{k}(1 - \bar{z}_{j}z_{k})^{-1}.
\end{equation*}
于是 $\|\tilde{B}h\|_{2} \leq \|h\|_{2}$ 的充要条件是,对任意 $\{\alpha_{j}\}_{j=1}^{n}$,
\begin{equation*}
\sum_{j,k=1}^{n} \frac{1 - w_{j}\bar{w}_{k}}{1 - z_{j}\bar{z}_{k}} \bar{\alpha}_{j}\alpha_{k} \geq 0.
\end{equation*}
\end{proof}

\begin{theorem}[Pick]\label{thm:2.17.7}
设 $\{z_{j}\}_{j=1}^{n} \subset D$ 且 $\{w_{j}\}_{j=1}^{n} \subset \mathbb{C}$.存在满足 $f(z_{j}) = w_{j}$($1 \leq j \leq n$)的 $f \in H^{\infty}$,且 $\|f\|_{\infty} \leq 1$ 的充要条件是
\begin{equation*}
\sum_{j,k=1}^{n} \frac{1 - w_{j}\bar{w}_{k}}{1 - z_{j}\bar{z}_{k}} \alpha_{j}\bar{\alpha}_{k} \geq 0 \quad (\{\alpha_{j}\}_{j=1}^{n} \subset \mathbb{C}).
\end{equation*}
\end{theorem}

\begin{proof}
若存在满足 $f(z_{j}) = w_{j}$($1 \leq j \leq n$)且 $\|f\|_{\infty} \leq 1$ 的 $f \in H^{\infty}$,则由引理~\ref{lem:2.17.3} 与 $\|T_{f}^{*}\| = \|T_{f}\|$ 和引理~\ref{lem:2.17.6},
\begin{equation*}
\sum_{j,k=1}^{n} \frac{1 - w_{j}\bar{w}_{k}}{1 - z_{j}\bar{z}_{k}} \alpha_{j}\bar{\alpha}_{k} \geq 0 \quad (\{\alpha_{j}\}_{j=1}^{n} \subset \mathbb{C}).
\end{equation*}
反之,若
\begin{equation*}
\sum_{j,k=1}^{n} \frac{1 - w_{j}\bar{w}_{k}}{1 - z_{j}\bar{z}_{k}} \alpha_{j}\bar{\alpha}_{k} \geq 0,
\end{equation*}
令 $B(1 - \bar{z}_{j}z)^{-1} = \bar{w}_{j}(1 - \bar{z}_{j}z)^{-1}$,由引理~\ref{lem:2.17.6},$\|\tilde{B}\| \leq 1$.因 $\tilde{B}T_{z}^{*} = T_{z}^{*}\tilde{B}$,令 $A = \tilde{B}^{*}$,则 $T_{z}A = AT_{z}$.于是由定理~\ref{thm:2.17.5},存在 $f \in H^{\infty}$,使得 $A = T_{f}$ 且 $\|A\| = \|f\|_{\infty}$.由引理~\ref{lem:2.17.3},$T_{f}^{*}(1 - \bar{z}_{j}z)^{-1} = \overline{f(z_{j})}(1 - \bar{z}_{j}z)^{-1}$,故 $f(z_{j}) = w_{j}$ 且 $\|f\|_{\infty} = \|\tilde{B}\| \leq 1$.
\end{proof}
